\scanpage{24}
\mspage{5}

\[
\varphi^{p'+q}\psi=\varphi^{p+q'}\psi
\]
\margin{On veut conclure $p'+q=p+q'$ i.e. $q-p=q'-p$}

On va supposer
\textcircled{B} Si $r,s\in\mathbb N$ sont tels que $\varphi^r\psi=\varphi^s\psi$, on a $r=s$.

On trouve alors une application
\diagram{\begin{tikzcd}[ampersand replacement=\&]M_0\arrow[r,"\delta"]\&\mathbb Z\quad\deleted{\unsure{g}},\end{tikzcd}}
et il est immédiat que celle-ci est \unsure{multiplicative|multiplication} ; i.e. $\delta(fg)=\delta(f)+\delta(g)$ :
\[
\left\{\begin{aligned}\varphi^p\psi f&=\varphi^q\psi\\ \varphi^{p'}\psi g&=\varphi^{q'}\psi\end{aligned}\right.
\]
\margin{multiplier par $g$ à droite}
implique $\varphi^p\psi fg=\varphi^q\psi g$
et en multipliant par $\varphi^{p'}$ à gauche
\[
\varphi^{p+p'}\psi fg=\varphi^{q+p'}\psi g=\varphi^q\bigl(\underbrace{\varphi^{p'}\psi g}_{\varphi^{q'}\psi}\bigr)=\varphi^{q+q'}\psi
\]

donc $\delta(fg)=(q+q')-(p+p')=\delta(f)\unsure{+|-}\delta(g)$.

Donc $\delta$ se factorise en
\diagram{\begin{tikzcd}[ampersand replacement=\&,column sep=large]M_0\arrow[r,"\mathrm{can}"]\arrow[rr,bend right=43,"\delta"']\&\underset{\substack{\text{groupe}\\\text{enveloppant}\\\text{de }M_0}}{G_0}\arrow[r,"\delta_0"]\&\mathbb Z.\end{tikzcd}}

% archive sheet number: 23
